
%%%%%%%%%%%%%%%%%%%%%%%%%%%%%%%% SCRIPT FOR E-RARA AND MDZ FILES     %%%%%%%%%%%%%%%%%%%%%%%%%%%%%%%%%%%%%%%%%%%%%%%%
%%%%%%%%%%%%%%%%%%%%%%%%%%% fini le 30.04.2025 par F. GOY            %%%%%%%%%%%%%%%%%%%%%%%%%%%%%%%%%%%%%%%%%%%%%%%%
% !TeX TS-program = lualatex
\documentclass{article}
\usepackage[T1]{fontenc}
\usepackage{microtype}
\usepackage[pdfusetitle,hidelinks]{hyperref}

\usepackage{polyglossia}
\setmainlanguage{english}
\setotherlanguages{latin,greek}
\usepackage[series={},nocritical,noend,noeledsec,nofamiliar,noledgroup]{reledmac}
\usepackage{reledpar}

\usepackage{fontspec}
\setmainfont{TeX Gyre Termes}

\usepackage{sectsty}
\usepackage{xcolor}

\usepackage{fancyhdr}
\pagestyle{fancy}
\fancyhf{}
\fancyhead[LE,RO]{\nouppercase{\leftmark}}  
\cfoot{\thepage}
\renewcommand{\headrulewidth}{0.4pt}

% Redefine \section to remove numbering
\usepackage{titlesec}
\titleformat{\section}[block]{\normalfont\scshape\color{gray}}{}{0pt}{} % no number in heading
\titleformat{\subsection}[hang]{\normalfont}{}{0pt}{} % also remove subsection number
\titleformat{\subsubsection}[hang]{\normalfont\footnotesize\color{black}}{}{0pt}{}

% Modify how section marks are stored to exclude numbers
\makeatletter
\renewcommand{\sectionmark}[1]{%
	\markboth{#1}{}} % Only store the section title, without number
\renewcommand{\subsectionmark}[1]{%
	\markright{#1}} % Only store the subsection title, without number
\renewcommand{\numberline}[1]{} % Hide the section number in TOC
\makeatother

\begin{document}

\date{}
        \title{Annotationes in Epistolas Pavli Ad Timotheum : [Bugenhagen, Johannes], [1524]}
\maketitle
\tableofcontents
\clearpage
\begin{pages} 
\beginnumbering
        ❧ IN PRIO- REM AD TIMOTHEVM EPI- stolam Annotationes Ioannis Bugen- hagii Pomerani. 
\phantomsection
\addcontentsline{toc}{subsection}{\textit{ARGVMENTVM. }}
\subsection*{\textit{ARGVMENTVM. }}\pstart \huge\textbf{T}\normalsize IMOTHEVM APOSTO- lum, id est, episcopum siue praedica- torem, a se ad hoc ordinatum Paulus instruit quid doceat, nempe fidem tan- tum et charitatem, reiectis fabulis et inutilibus quae stionibus, et pugnis uerborum, et di- Jputationibus inanibus pro ambitione et auaricia ina stitutis, quae pestes sunt syncerae in Christum fidei, et in proximum charitatis: neue admittat doctrinas dae moniorum in hypocrisi loquentium mendacium, quas in nouissimis temporibus regnaturos praedicit. Haec spar- Iim in epistola, cap.2. de oratione communi instruit, et quid mulieres deceat erga uiros. In tertio, Episcopos, et Diaconos Christianos depingit cum corum uxori- bus, familia et liberis. In quinto modum monendismn- gulos praescribit, et uere uiduis prouideri uult ab ecclesia. Docet quoque Episcopos, qui laborent in- uerbo et doctrina, esse dignos, quibus uictu et ami- ctu prouideatur, et de iudicio episcopi in peccantes lubdit. In sexto docet, quod  serui, quoddiuites monendi sint?  \pend
\section*{ANNOT. IOANNIS POMERANI }\pstart praeterea, quae de doctrina Christiana interserit:ut su- pra dictum est.  \pend
\phantomsection
\addcontentsline{toc}{subsection}{\textit{LOCI INSIGNIORES SVNT. }}
\subsection*{\textit{LOCI INSIGNIORES SVNT. }}\pstart Primus, quae sit doctrina Christiana, et quae pro- phana, locus maxime necessarius, ne fiducia operum legis, iustificatorem deum negemus, aut inani disputa- tione et pugna uerborum, et inutilibus quaestionibus. fidem simul laedamus et charitatem, cap.1.4.et 6. Se- cundus, quod lex iusto non est posita, sed iniustis, cap. I.id quod ad Gal.5. sic legis: Si spiritu ducimini, non estis sub lege: manifesta sunt autem opera, etc. Terti us, quod blaspemia et contumelia et persecutio con- tra Euangelium Christi, non noceant auctoribus suis, qui trrantes per hoc putant se obsequium praestare deo modo postea cognita ueritate resipiscant, ut Pau. Alii uero qui scientes contra ueritatem pugnant, et blas- phemant, peccant irremissibiliter inspiritum sanctum, ut est in euangelio. Verum hic non est nobis praecipitan- da sententia. Hic locus hodie scitu necessarius est. cap. 1. Quartus, quod pro omnibus est orandum. item pro regibus et praelatis, ut sub eis nobis quiete uiue- re liceat. Deus enim nullum statum reiicit: sed ut ex iu- fimis, ita et ex sublimibus quosdam ad sui agnitionem perducit, cap.2. Quintus, habitus mulierum, et ut sint in omnibus uiris subiectae, cap.2. Sextus, officium et shectata conuersatio Episcoporum et Diaconorum,  \pend
\section*{IN EPI. PAV. AD TIM. I. }
\marginpar{[ p.79 ]}\pstart et eorum uxoru, liberorum et familiae, cap.3. Septi- mus, doctrina daemoniorum. Satis hodie notus locus, cap.4. Octauus, pietatem ad omnia ualere, parum uero torporis exercitationem, locus magis necessarius, quam uideatur, cap, 4. Nonus, de uiduis, cap, 5. Deci- mus, quod necessaria uitae debentur doctoribus ab ipsis auditoribus, cap. 6.Vndecimus, de accusatione et iu- dicio in peccantes, cap.5. Duodecimus, quod neque  ser- ui, neque  diuites alieni sunt a regno coelorum, etiam in ipsa seruitute, et in ipsis diuitiis, modo sequantur quae Paulus monet, cap.6. Praeterea qua comitate agendum lit cum omnibus, principium cap. 5. ostendit.  \pend
\endnumbering\beginnumbering\section{CAPVT I.}
\phantomsection
\addcontentsline{toc}{subsection}{\textit{Paulus apostolus Iesu Christi. }}
\subsection*{\textit{Paulus apostolus Iesu Christi. }}\pstart \huge\textbf{H}\normalsize Ac inscriptione suum uerbum commendat Ti- motheo, et omnibus lecturis epistolam, ut certi- simus hoc esse uerbum dei, non hominis: quandoquidem do- ctrinas et mandata hominum reiicit deus in Esaia, et Christus in euang. et lex non patitur aliquid addi uera- bo dei. Dicit ergo se apostolum esse Christi, et sibi com- missum esse apostolatum, id est officium euangelii prae- dicandi a deo. quasi diceret, non uenio mea referre, sed quae mihi commissa sunt a deo.  \pend\pstart Deum saluatorem uocat, quia ipse per Christum seruat, siue saluat nos: et ficut dicitur de Ephe.1.fecit nos di-  \pend
\section*{ANNOT. IOANNIS POMERANI }\pstart lectos in dilecto filiosuo. Et Christu uocat spem no- stram, quia qui Christum non habent, in desperatione uiuunt, sicut ad Ephe.4. Qui desperant, semetipsos tra diderunt in impudicitiam etc. Hypocritae quoque  fin- gunt se sperare, sed ueniente tentatione, spem non ha- bent, quia Christum non habent. Germanum, id est uerum- nam Christum recte didicerat ex doctrina Pauli, qui se ob id eius patrem scribit. Qui germani filii in fide sint, parabola de seminante, Lucae. 8. declarat.  \pend\pstart Gratiam et miserocordiam (quae est, quod gratis no- bis donantur peccata, et haeredes simus regni absque  omni merito) necessario sequitur pax mentis et consci entiarum, quae exuperat omnem mentem. Non enim oculus uidit, nec auris audiuit, etc. Nisi uero gratiam dei, et misericordiam te consecutum credideris, nunquam pacem habebis. Consecutum inquam, non ex tuis meri- tis, iustitiae studio, etc. sed a deo, qui tantum nos dili- git in dilecto filio suo, ut pater noster non solum ap- pelletur, sed reuerasit: ut nihil trepidemus ad tantam maiestatem, qui ex misericordia eius sumus filii. Et a do- mino Iesu Christo domino nostro, sub cuius dominio et imperio totus exercitus inferorum, dei filiis, et Chri- sti seruis praeualere non potest, etc. Atque hic uides Christum non solum dominum dici, ne quis contemnat eum, quia uidet hominem factum: sed etiam dominum nostrum uocat, ut etiam agnoscas eum protectorem,  \pend
\section*{IN EPI. PAV. AD TIM. I. }
\marginpar{[ p.80 ]}\pstart imo et te dominumi, et regem in ipso, qui potens sis lupra peccatum, mortem, et infernum etc.  \pend
\phantomsection
\addcontentsline{toc}{subsection}{\textit{Ne diuersam, etc. }}
\subsection*{\textit{Ne diuersam, etc. }}\pstart Vel ne aliam sequantur doctrinam, ut uetus habet translatio. Ne aliter doceret, praedicationi euangeli- cae siue fidei nihil uult addi, quantacumque  specie pieta- tis: id quod tum quoque  nitebantur pseudapostoli. Vnde uit ad Gal. 1. Etiam si angelus e coelis etc. Nimirum prae- sentiebat fidem Christi, doctrinis humanis abolendam, ut factum cernimus. Idcirco summe notandum, quod hic dicit, aedificationem dei esse perfidem, sentiens re- liqua omnia, quae docentur, esse destructionem. Aedi- ficatio est doctrina salutis: destructio, impia doctrina, et salutis uastatio, ut saepe uides in Paulo. Hanc fidei aedificationem duo hominum genera impediunt, uel etiam destruunt aedificata. Nam alii docent fidere in operibus legis: quia lex a deo data, non potest non esse bona et necessaria. Quibus respondet, quod usum le- gis bonae ignorent. Fidem oportet doceri, sine qua le- gis opera, quemadmodum lex exigit, fieri non possunt. eides enim spiritum impetrat, legis impletorem: ubi uero habueris fidem, iam nulla lege teneris, ut nemi- ni quicquam debeas, nisi ut proximum diligas, ad Rhoma. 12. Et hoc praestabis, non solum quia debes, quia lex iubet: sed et quia uis, nimirum sic trahente spiritu uoluntatem tuam. Hoc est quod hic dicit, lex sint. iusto non est posita, sed iniustis, qui uel puniuntur, uel  \pend
\vspace{0.5cm}\noindent
\vspace{0.2cm}\rule{1cm}{0.2pt}\\ 
\hspace{0.2cm}\textit{mg}
\footnotesize An sontes plectendi 
\normalsize| 
\section*{ANNOT. IOANNIS POMERANI }\pstart occiduntur secundum legem:ut non sit opus quaerere, an sub nouo testamento sint sontes plectendi, quum constet adulteros, fornicarios, homicidas sub nouo testamento non contineri, ad quod pertinent soli iusti per fidem,is est credentes, contra quos non est lex. Deinde cum multa ex lege praecipiant, intentionem legis et capute egregie negligunt, quae est, ut dixi, charitas in proximum, sed non foris simulato opere ostenfa:sed, ut hic definitur, ex corde puro, et conscientia bona, et fide non ficta. Vides quod ne eleemosyna quidem quantuncunque co- piosa, prosit sine fide (quicquid enim non est ex fide peccatum est) et charitatis opera quae uidentur, non esse charitatem, nisi ex corde et conscientia, quam fides purificauit, non ficta ueniant. Ii, inquit, deflexerunt ad uaniloquium, uolentes esse legis doctores, sicut et hodie operum praedicatores et doctores, non intelli- gentes quae loquuntur, nec de quibus adseuerant. Prae cepta quidem proponunt, et legem siue mandata mo- rum: sed unde tibi detur, ut sic uiuas, non docent, Inue- nies hodie qui nihil quaerant in sacris literis praeter moralia, quae deinde similia putant ethnicis, Aristote- licis, aut Senecae praescriptis: qui tamen nesciunt, quid Paulus in tota epistola ad Rhoma. uel Gal. disputex aut quid Christus uelit, dum fidem, id est fiduciam in deum ubique  praedicat: aut quid sit quod hic Paulus dicit, cha- ritas de puro corde, et conscientia bona, et fide non  \pend
\section*{IN EPI. PAV. AD TIM. I. }
\marginpar{[ p.81 ]}\pstart ficta. Nam, ut inquit, quod a talibus aberrauerunt, deflexerunt ad uaniloquium, etc. Sed de his hactenus. Alii uero docti ex sacra scriptura, sciunt non fiden- dum operibus, sed soli dei misericordiae, et quidem recte: sed curiosi et superstitiosi inania uocabula in- ueniunt, quae literae sacrae non nouerunt, et scientiam iactant, contra quos in fine epistolae dicit, deuitans pro phanas uocum inanitates. Item fabulas sequuntur, id est, alia extra sacras literas: qualia multa commenti sunt Iudaei, qualia etiam multa nostri habent quasi hi- storiam. Et praeterea genealogias quaerunt numquam finiendas, id est de quibus ex sacris literis non potes esse certus, quas alius sic, alius aliter computat, et quisque  sua praefert: quae, inqt, quaestiones et dubium ge- nerant, et nullam aedificationem. Item sunt, quorum magnus est numerus, qui magno animi negotio, nihil nisi uanas quaestiones sequuntur ex sacris literis, nus- quam non haesitantes, dum etiam nodum inscirpo quae- runt: interim uero quam certissimam uident sanam doctrinam, et apertum Euangelium, id est luce clariores: dei Christique  promissiones in utriusque  instrumenti pa- gina relinquunt. Qui in hoc periculo sunt, ut ueren- dum sit eis, ne dum pugnas sequuntur uerborum, reli- ctis illis, ex quibus aedificari possent, tandem incipiant de ueritate sacrae scripturae dubitare. Et dum umbram sequuntur, ut canis Aesopicus, ueritatem et rem ipsam:  \pend
\section*{ANNOT. IOANNIS POMERANI }\pstart id est, fidem amittant, quorum adhuc peior est causa, si per haec sua lucra sectantes, alii aliis sese praeferre conten- dant. de quibus cap. 6. sic dicit: Si qs aliter docet, etc.  \pend
\phantomsection
\addcontentsline{toc}{subsection}{\textit{Ex corde puro. }}
\subsection*{\textit{Ex corde puro. }}\pstart Vbi enim est fides non ficta, ibi est et cor purum: si- cut Petrus dicit: Fide purificans corda eorum. ibi est et conscientia bona. Qui enim sibi recte conscius est qui deo non ex toto corde fidit? Porro fidem non fictam dixit, id est ueram: contra illam quae hodie uulgo iacta tur, qua adsentiuntur historiae euangelicae: sicut aliis historiis, et illis quae gesta narrantur de Thurca, de Alexandro, de Iulio papa, etc. quae potius opinio est de Christo, et fides historica, uel potius ficta fides quam fides illa iustificans, quamsoli norunt credentes, Hinc fit, ut nostri libentius audiant fabulam narrari quae exempla uocant, quam textum euangelicum: quia quam fidem Euangelio habent, etiam habent illis fabu- lis non aliam, et fere minorem.  \pend
\phantomsection
\addcontentsline{toc}{subsection}{\textit{Et gratiam habeo. }}
\subsection*{\textit{Et gratiam habeo. }}\pstart Hic locus nobis proponitur a Paulo, ne ob quan- libet etiam uitam, et quaecumque  peccata desperemus. modo ueniamus ad Christum, ut Paulus. Hic uide quam sint omnia dei, ut hic nihil praesumas, sicut illi qui stul- te putant sese posse ad Christum aditum habere quando- cunque  uolunt, contra Paulit, dicentem: Non est uolentis,  \pend
\vspace{0.5cm}\noindent
\vspace{0.2cm}\rule{1cm}{0.2pt}\\ 
\hspace{0.2cm}\textit{mg}
\footnotesize Rhom.9 
\normalsize| 
\section*{IN EPI. PAV. AD TIM. I. }
\marginpar{[ p.82 ]}\pstart neque  currentis, sed etc. Primum dicit, qui me potentem, reddidit per Christum: deinde, qui me fidelem iudicauit. item, misericordiam adeptus: item, exuberauit supra mo- dum gratia etc. item: Christus uenit peccatores saluos. facere, etc. Non nobis domine non nobis, sed nomini tuo da gloriam. Super misericordia tua et ueritate tua. id quod Paulus hic sic dicit: Regi autem seculorum, etc.  \pend
\phantomsection
\addcontentsline{toc}{subsection}{\textit{immortali regi. }}
\subsection*{\textit{immortali regi. }}\pstart Ergo regnum eius est in aeternum. Inuisibile ergo re- gnum eius non est de hoc mundo, sed spirituale. Solisa- pienti, ergo praeter eum et extra eum omnia sunt stula- ticiae plena: ut nihil sapientiae uerae sit, ubi ipse non il- luminat: et, ut sic dicam, sapientificat. Et de Christo praedictum erat, Hiere. 23. Et regnabit rex, et sapiens erit. Ex Hebus  sanctificator erit.  \pend
\phantomsection
\addcontentsline{toc}{subsection}{\textit{Iuxta prophetias. }}
\subsection*{\textit{Iuxta prophetias. }}\pstart Intelligunt, iuxta ea, quae spiritu prophetiae ante de hoc cognoui, quando te feci episcopum, ut secundum ea a- gas. Ego aunt et hic, et infra cap.4. intelligo prophe- tias doctrinas spiritus dei, siue reuelationem scripturae, quam a Paulo uel aliter acceperat Timotheus: sicut prophetiam accipi diximus in Esaia. Vnde ex graeco hic le- gi potest iuxta prophetias, quae ad te praecesserunt: id est ut ambules secundum doctrinam quam accepisti, et quam habes, ut eam ipsam defendas contra aliter docentes:  \pend
\vspace{0.5cm}\noindent
\vspace{0.2cm}\rule{1cm}{0.2pt}\\ 
\hspace{0.2cm}\textit{mg}
\footnotesize Psal.113 
\normalsize| 
\hspace{0.2cm}\textit{mg}
\footnotesize vide Anno tationes e- iusdem in E- saiam. 
\normalsize| 
\section*{ANNOT. IOANNIS POMERANI }\pstart idquefacias in fide, et bona coscientia, quae duo separa- ri non possunt. Nam, ut inquit, qui conscientiam bonam repulerunt, id est qui audent aliud dicere, aliudquefa- cere, quam quod in corde nouerunt, unde aliquid prae caeteris uideantur, unde caeteris superiores per conten- tionem adpareant: id quod hodie multum uulgare est. Sed Paulus hic uocat blasphemiam. Hi, inquit, qui bo- nam conscientiam repulerunt, naufragium fecerunt circa fidem: id est, etiam in eo periclitati sunt, quod intus in corde sciuerunt, et tandem perdiderunt.  \pend
\phantomsection
\addcontentsline{toc}{subsection}{\textit{Hoc praeceptum }}
\subsection*{\textit{Hoc praeceptum }}\pstart Scilicet, ut supra, ut denuncies quibusdam etc. ut com- mendes gratiam Christi peccatoribus, sicut ego, absque  operibus legis, etc. Vel potes referre ad hoc quodse- quitur:ut milites, etc.  \pend
\phantomsection
\addcontentsline{toc}{subsection}{\textit{Hymenaeus et Alexander. }}
\subsection*{\textit{Hymenaeus et Alexander. }}\pstart Qui scientes, contra ueritatem quam sciuerunt, so- Ient contendere.  \pend
\phantomsection
\addcontentsline{toc}{subsection}{\textit{Tradidi satanae. }}
\subsection*{\textit{Tradidi satanae. }}\pstart Sicut Corinthium, qui duxerat nouercam suam: de quo uide in annotationibus Philippi. Hanc tamen ex- communicationem, sicut et illam, ostendit esse utilem, cum dicit: Vt discant non blasphemare. Haec non est pa- pistica, quam hodie uidemus excommunicatio, ett,  \pend
\section*{IN EPI. PAV. AD TIM. I. }
\marginpar{[ p.83 ]}
\endnumbering\beginnumbering\section{CAPVT II.}\pstart \huge\textbf{O}\normalsize Ratio est cordis desyderium pro re a deo im- petranda. Hoc desyderium si uerum est, nun- quam cessat, donec a deo quod desyderat, accipit: hoc est quod Christus dicit, Oportet semper orare, et non deficere. Qui non desyderat, nihil orat, ut maxime multiloquio aëra compleat: qui uero desyderat, etiam a uerbis quandoque  non abstinebit. Hinc fit, ut generali uocabulo orationem uocemus quamcunque  cum deo col- locutionem, etiam cum laudamus in psalmis et canti cis spiritualibus: id est, quae fiunt incordibus. Nam ut nunquam cessamus desyderare misericordiam et no- bis et proximis: ita quoque  semper, ne nostra uidea- mur quaerere, desyderamus sanctificari nomen patris nostri. In omnibus ad deum fiat uoluntas tua etc. In- struit ergo hic Paulus episcopum, quod moneat suos auditores, nempe hoc, ut in primis orent. Quomodo enim intelligent Euangelium, aut dilectionem conse- quentur in deum et in proximum, nisi postulauerint a deo spiritum bonum? Vt hac admonitione primum omnium discant omnia ex deo pendere, cum audiunt omnia a deo postulanda. Orent autem pro omnibus hominibus, etiam inimicis: id quod uera postulat cha- ritas. Nec excludant ab oratione reges et principes, licet impios. Primum, ut quam pacem eis postulamus ueniat quoque  ad nos, et sub eis uiuamus in pace, quan-  \pend
\vspace{0.5cm}\noindent
\vspace{0.2cm}\rule{1cm}{0.2pt}\\ 
\hspace{0.2cm}\textit{mg}
\footnotesize Oratio quod . 
\normalsize| 
\section*{ANNOT. IOANNIS POMERANI }\pstart tum per pietatem ad deum, et honestatem ad proxi- mum liceat. Ne intelligas propter carnis pacem, quod ad homines attinet, cedendum esse pietati et fidei, si- ue conscientiae. Deinde, quod hoc quoque  acceptum est deo, ut pro omnibus oremus sine acceptione persona- rum, siue sint iudaei siue gentes, siue serui siue |principes. Quia sicut ipse solus deus est, et omnium sine acceptio- ne dominus: ita oens homines uult ad agnitionem sui per- uenire. Quapropter nos neminem excipere debemus, ne reges quidem impios, cum ipse neminem excipiat. Sic pratur pro Nabuchodonosore, et Balthasaro, Baruch. 1.Item, quia unus est mediator inter deum et hominem, qui se dedit redemptorem pro omnibus. Ergo rursum nos neminem excludere debemus, cum oramus pro sa- lute hominum, etc. scientes hoc ipsum et deo patri, et Christo, quem pater dedit mediatorem, placere.  \pend
\phantomsection
\addcontentsline{toc}{subsection}{\textit{Deprecationes. }}
\subsection*{\textit{Deprecationes. }}\pstart Sic distinguunt, ut deprecari sit orare, ut malum aufe- ratur, obsecrare ut bonum ueniat, interpellari, conque- ri de his quae nos laedunt: gratias agere de beneficiis acceptis: id quod fere rationes uocabulorum postulant.  \pend
\phantomsection
\addcontentsline{toc}{subsection}{\textit{Omnes uult saluos fieri. }}
\subsection*{\textit{Omnes uult saluos fieri. }}\pstart Cunctos homines uult saluos fieri, siue seruari, et ad agnitionem ueritatis peruenire: ergo iussi sunt aposto- li praedicare Euangelium omni creaturae: Sed multi sunt uocati, pauci uero electi Quomodo autem hoc con-  \pend
\section*{IN EPI. PAV. AD TIM. I. }
\marginpar{[ p.84 ]}\pstart uenit cum eo quod Rho. 9. dicit Paulus, Cuius uult mise- retur, et quem uult indurat? Respondeo, quod unus lo- cus in scripturis obscurus, non est omnibus aliis claris obiiciendus. Tota scriptura habet impios damnari, pios siue credentes saluari: illos iudicio dei, hos miseri, cordia. Intelligendum ergo quod dicit, Vult ommes homines saluos fieri etc. De hominibus omnium statu- um, id quod hic cohaerentia textus postulat. Orate pro omnibus, etiam pro regibus etc. quia hoc placet deo: qui nullum statum mundi reiicit: quandoquidem et iu- stos reges fuisse constat:ut Dauidem, Ezechiam, etc. Omnes ergo hoïes, id est ex omnibus statibus homi- nes, et praeterea non solum Iudaeos, ut ipsi Iudaei putant: sed et gentes per totum orbem. Sic et de Christo dicit: Qui dedit semetipsum precium redemptionis pro omni- bus. Vide, si placet, annotationes. Philippi ad Rho.8.  \pend
\phantomsection
\addcontentsline{toc}{subsection}{\textit{Vnus mediator. }}
\subsection*{\textit{Vnus mediator. }}\pstart Vnde ergo nobis tot mediatores quidam fecerunt Iine scriptura, et sine uerbo dei? quibus non occur- rit, quod Christus orat pro nobis, et intercedit, ut mediator: et solus mediator inter deum et hominem, quam solus est deus, et homo: id quod nullus aliorum sanctos- rum esse potest. Orant sancti qui sunt in terris pro se mutuo, et exaudiuntur a deo, sed per Christum media- torem, sicut ipse promittit: Si quid petieritis patrem in nomine meo, hoc faciet, etc. De defunctis sanctis,  \pend
\vspace{0.5cm}\noindent
\vspace{0.2cm}\rule{1cm}{0.2pt}\\ 
\hspace{0.2cm}\textit{mg}
\footnotesize Ioan.16. 
\normalsize| 
\section*{ANNOT. IOANNIS POMERANI }\pstart quod orent pro te, scripturam et uerbum dei non habes. Quan- do quidem scriptura Christum tibi commendat pontificem qui ingressus est in coelum, ut illic se offerat uultui patris pro nobis, ut dignissime tractat epistola ad Hebus  Qui in- terpellat pronobis, Rho.8. Quem aduocatum habemus apud patrem, qui est propiciatio pro peccatis nostris, I. Ioan.2. Et propiciatorium nostrum, ad Rho.3. Et thronus gratiae, ad Hebus 4. Et praeterea solum deum inuocandum esse tota scriptura et praecipit, et testa- tur, nec aliud exemplum inscripturis habes. Cur ergo ad illa deflectis quae nescis, relicto illo quod certum ti- bi offertur? In iniuriam certe Christi alios in coelo po- suerunt mediatores. Et Paulus aperte dicit, vnus de- us unus mediator.  \pend
\phantomsection
\addcontentsline{toc}{subsection}{\textit{Vt esset testimonium temporibus suis. }}
\subsection*{\textit{Vt esset testimonium temporibus suis. }}\pstart Id est, ut hoc mysterium redemptiomis Christi, qua redimuntur omnes homines: id est non solum Iudei, sed etiam gentes, reuelaretur mundo, tempore a deo ordi- nato: sicut de hoc mysterio reuelando in epistola ad Ephe. et alibi scribit, et hic infra cap.3. Ad quod my- sterium reuelandum dicit Paulus se doctorem missum à deo gentibus. Doctorem inquam in fide, non doctrinis le- gis, et ueritate, non hypocrisi operum. Veritas est fides, ubi creditur non hominum doctrinis (omnis enim homo mendax) sed soli uerbo dei: qui est ipsa ueritas, etc.  \pend
\phantomsection
\addcontentsline{toc}{subsection}{\textit{Volo uiros orare. }}
\subsection*{\textit{Volo uiros orare. }}
\section*{IN EPI. PAV. AD TIM. I. }
\marginpar{[ p.85 ]}\pstart Prosequitur quod oeperat de oratione. Vide hic quod Paulus non requirat templum adorationi: sed dicit, In omni loco, id est, ubicunque  orauerint, etc. Ne hic stru- as rursum Paulo calumniam, quod oporteat etiam in aquis, in inuiis etc. orare: quia dicit in omni loco, etc. Potueram et hoc quoque  quod supra dixi, Vult omnes homines saluos fieri: sic interpretari, id est, quotquot saluantur, eius uoluntate saluantur, sicut Ioannis I. Illu- minans omnem hominem etc. id est quotquot illuminan- tur, a deo illuminantur: sicut et hic, Volo uiros orare in omni loco:id est, ubicunque  orauerint. Si uidissem ad contextum orationis pertinere, ut uideas uel hoc loco, quod scriptura nihili facit sophisticas cauillationes. Qui credit, credat. Argutiis enim illis, deispiritus non conci- pitur. Summa ergo (ut hoc rursum ingeram) superioris loci ex contextu orationis petenda est. Orate pro omni- um hominum statu siue conditione. Nam deus nullum statum reiicit, et Christus non solum pro Iudaeis, sed etiam pro gentibus, cuiuscunque  conditionis existant, mortuus est. Aduerte, quod uiros uult orare cum pietate, mulieres ue- ro non solum cum pietate, sed et cum uestitu et opere profitentes pietatem, ut etiam ipso aspectu uideantur esse Christianae: et exemplumsint castitatis et uere- cundiae, ne auertant uirorum animos, ubi ad orationem publicam conuenitur. Nam publicam etiam orationem tum habuisse Christianos, ubi conueniebant, testatur  \pend
\section*{ANNOT. IOANNIS POMERANI }\pstart Iocus Pauli 1.Corin. 14. Si benedixeris spiritu, qui supplet locum idiotae, quomodo dicet amen super tuam benedictionem. Opera uero bona mulierum non sunt fuperstitiones illae stultarum foeminarum, quarum mul- tas ipsae excogitarunt, multas uero a stultis confesso- ribus, quos uocant, et indoctis praedicatoribus didi- cerunt. Sed sunt illa quae scribuntur infra cap.5. Si filios educauerit, si fuit hospitalis, si sanctorum pedes lauerit, si adflictis inseruiuerit, si in omni bono opere fuerit adsidua. Et ad Titum 2. Anus similiter et c.1 Petri 3. Similiter et mulieres, etc.  \pend
\phantomsection
\addcontentsline{toc}{subsection}{\textit{Sustollentes etc. }}
\subsection*{\textit{Sustollentes etc. }}\pstart Habitum orantium indicat, qui nihil est nisi et cor le- uetur ad deum: sicut in psalmo dicitur: Ad te domine leuaui animam meam: id quod hic dicit, Sustollentes pu- ras manus. Quorum enim manus sanguine plenae sunt, illos non uult exaudire dominus, Esa.1. Purae manus non sunt, nisi cor purumsit, quod sine fide purum esse non potest. Fide enim, inquit Petrus, purificans corda eorum, id quod hic dicit, absque  ira et disceptatione. Duo sunt quae impediunt exauditionem. Primum ira, id est quod non dimittis proximo in te peccanti. Nam hanc conditionem in euangelio semper uides adiectam: Si re- miseris fratri, etc. Alterum disceptatio, qua disceptas tecum, et dubitas num sis exaudiendus. Hic quoque  nihil re- cipis, Iacobi. I. Vide Marci II. Credite et accipietis.  \pend
\vspace{0.5cm}\noindent
\vspace{0.2cm}\rule{1cm}{0.2pt}\\ 
\hspace{0.2cm}\textit{mg}
\footnotesize Bona opera mulierum. 
\normalsize| 
\hspace{0.2cm}\textit{mg}
\footnotesize Duo sunt quae impe- diunt exau- ditionem. 
\normalsize| 
\section*{IN EPI. PAV. AD TIM. I. }
\marginpar{[ p.86 ]}
\phantomsection
\addcontentsline{toc}{subsection}{\textit{Mulier in silentio, etc. }}
\subsection*{\textit{Mulier in silentio, etc. }}\pstart Id est, non loquatur, ubi docendi gratia conuenitur, ubi audiendum est uerbum dei:nam uiris permittitur Ioqui cum doctore, imo si fuerit assidenti reuelatum, qui prius docebat, taceat, 1. Corin.14. Mulieres uero in ecclesia taceant, ut dicitur ibidem: ne sibi authorita- tem uendicent in uiros contra legem dei, quae dicit: Sub uiri potestate eris. Et praeterea, Adam est creatione prior. Item, Mulier est ruina uiri: salua tamen fiet per liberorum generationem. Quod opus haud dubie bonum, nemo hodie respicit, cum tamen hic mulier habeat uer- bum Pauli, imo dei in Paulo, et in hoc obsequitur deo, qui dicit: In labore paries, etc. Imo hic in dei opere est, qui dicit: Multiplicabo conceptus tuos, et aeru- mnas tuas, etc. Et additur: Si manserint, etc. Sic enim est in Graeco uerbum plurale, quod in ueteri transla- tione non inepte legitur in singulari, uitato Hebraismo. Hebraismus enim est, si manserint scilicet mulieres. de quibus supra dixit. Consimiliter et mulieres, etc. Hebraeos subito mutare numerum, et personas, ex prophetis notum est. Nam qui intelligunt, si perman- serint, scilicet filii, non uident quam ineptum hoc sit, quasi mulier non possit salua fieri, nisi et filii saluentur. Igitur si permanserint, inquit, mulieres in fide, et dilectione et sanctificatione, qua quotidie sanctiores fiunt, id est trescunt fide et dilectione, adiuncta castitate, id est morum integritate, quae etiam deceat coram hominibus. Aduer-  \pend
\section*{ANNOT. IOANNIS POMERANI }\pstart e uero et mulieres prophetare, si habent uerbum, ubi uiri non sunt qui uerbum habeant. I. ad Corin.11. Mulier orans aut prophetans, etc. Iohelis 2.legis, Et prophetabunt filii uestri, et filiae uestrae. Actorum. 21. Erant Philippo quatuor filiae uirgines prophetantes. Et Maria canit Magnificat, etc.  \pend
\endnumbering\beginnumbering\section{CAPVT III.}
\phantomsection
\addcontentsline{toc}{subsection}{\textit{DE EPISCOPIS ET DIACONIS. }}
\subsection*{\textit{DE EPISCOPIS ET DIACONIS. }}\pstart \huge\textbf{E}\normalsize Piscopos esse uerbi dei praedicatores ad hoc ele- ctos, et diaconos esse sanctorum ministros, et pauperum prouisores admonui in principio epistolae ad Philip. Hi eligebantur in ciuitate a populo, uel ab Epi- scopo: id est, apostolo aliquo, siue praedicatore, qui illic populum docuerat: et manere ibi non potuit consentiem te tum et uolente populo. Eligebantur autem ex ciui- bus probatissimis, quibus erant uxor, filii, familia, domus cura, et. Sicut nunc consules eliguntur: ut hic uides. Vnde episcopatus et diaconatus officia erant, non perpetuae dignitates, ut nunc fabulantur de chara- ctere indelebili. Compone ergo nunc nostrorum dignita- tes illorum officiis, et uidebis apostolicam institutionen nostros ignorare. Igitur apostolus dat, si qui sint qui cupiunt praefici, quemadmodum hic describit. Bonum enim opus desyderant, non suam gloriam aut commodum: officia enim sunt, non dignitates. Eadem pene uides ad Titum  \pend
\vspace{0.5cm}\noindent
\vspace{0.2cm}\rule{1cm}{0.2pt}\\ 
\hspace{0.2cm}\textit{mg}
\footnotesize episcopatus officium, non dignitas est 
\normalsize| 
\section*{IN EPI. PAV. AD TIM. I. }\pstart I. quae huc quoque confer. Tales describit, qui non jo- lum docti sint et potentes in uerbo et fide euangelii Christi, et inculpati coram deo: sed etiam contra quos ne infideles quidem habeant, quod uel iuste obiiciant, aut saltem de quo male suspicentur, non solum ex pros pria persona: sed etiam ex uxore, filiis, familia, etc. ne blasphemetur bonum dei uerbum, et.  \pend
\phantomsection
\addcontentsline{toc}{subsection}{\textit{Vnius uxoris maritum. }}
\subsection*{\textit{Vnius uxoris maritum. }}\pstart Id est, qui non habeat simul plures uxores, quemad- modum nunc: quidam ex Iudaeis, quidam et ex genti- bus nuper conuersis adhuc habebant. Hoc ideo, ne sint lufpecti de impudicitia. Non ergo hic dicitur, quod mor tua uxore non possit ducere aliam, quasi tunc debeat resignare familiae, liberis, domui: quae hicuides, etc.  \pend
\phantomsection
\addcontentsline{toc}{subsection}{\textit{Nouicium. }}
\subsection*{\textit{Nouicium. }}\pstart Id est, nuper conuersum, et nondum probatum tene tationibus: quales fere solent putare se aliquid esse, cum nihil sint. Superbiunt enim quasi fide fortes, et ueniente tentatione cadunt, et negant, aperientes ora blasphemiis contra euangelium dei.  \pend
\phantomsection
\addcontentsline{toc}{subsection}{\textit{Tenentes mysterium fidei. }}
\subsection*{\textit{Tenentes mysterium fidei. }}\pstart Id est, cognoscentes et habentes apud se fidem, quae reuelatur in corde, quae a mundo non uidetur. Et hoc cum pura conscientia, sicut antea monuimus: fidem cĩt conscientia bona coniungi: id est qui boni non audeane  \pend
\section*{ANNOT. IOANNIS POMERANI }\pstart aliud dicere aut facere, quam quod secundum fiden sentiunt: hoc est quod in Actis apostolorum eligumtuy diaconi siue ministri, pleni spiritu sancto. sicut dixiin epistola ad Philippen.  \pend
\phantomsection
\addcontentsline{toc}{subsection}{\textit{Nam qui bene ministrauerint. }}
\subsection*{\textit{Nam qui bene ministrauerint. }}\pstart Id est cum pura confetentia, et abique crimine, si- cut supra dixit, ministrent sic, ut nemo possit illos cri- minari. Hoc uero quod hic additur, gradum bonum, et- uide ne intellexeris de gradibus ecclesiasticis, ut ho- die dicunt: quasi qui bonus fuerit aliquam diu diaconus possit deinde promoueri in episcopum, quod dignies sit maioribus, qui minora bene praestiterit, scilicet qui aliquandiu bonus fuerit fossor, potest deinde in bons- promoueri sutorem. Ita errant nostri, solum dignita- tes et mundanos honores somniantes, non officia:sci licet qui bene ministrauerit, et bonus minister fuerit ex dono spiritus, potest etiam ab hominibus promoueri ut sit bonus doctor plebis. Id enim est episcopum esse. Quid si donum prophetiae non habeat? ut maxime habe at mysterium fidei cum pura conscientia, sicut supra dixit quemadmodum omnes uere Christiani debent habere Spiritus nanque diuidit singulis prout uult. Tu ne illum- ex bono ministro promouebis in bonum episcopum? Simpliciter ergo hic uide, quod Christus dicit: Omni habenti dabitur, et abundabit, qui accipit donum a deo et utitur eo ad utilitatem ecelesiae, sicut Paulus dicit;  \pend
\section*{IN EPI. PAV. AD TIM. I. }
\marginpar{[ p.88 ]}\pstart Vnicuique datur manifestatio spiritus ad utilitate eccles siae, hic inquam habet, et huic eiusdem dono spiritus augetur illa gratia, et abundabit, ut possit adhuc ampli us praestare et prodesse. Hoc est quod hic dicit: qui bene ministrauerint, id est qui accepto hoc dono spiri- tus, ut possint esse ministri electi sanctorum, bene usi fuerint ad utilitatem ecclesiae, siue indigentium: gradum sibi bonum adquirunt, id est ex dono dei augentis et ditantis proficiunt, et altius ascendunt in eodem dono, ut indies magis ac magis quam ante pdesse possint Quia omni habenti dabitur, et abundabit. Nota est parabo la illa euangelica, quae certe de donis loquitur, ubi hoc additur: Ab eo autem qui non habet, id est dono accepto, sicut illic uides, non utitur ad sanctam usuram: etiam quod uide tur habere, auferetur ab eo, id est inutilis reddetur ab lato rursus dono, et paulatim ab eo etiam fides cordis aufertur, et excaecatur: quia infidelis fuit in commisso. Idcirco hic etiam (quod mihi maxime et unice placet) poteris coniungere sic: Gradum sibi bonum adquirunt, et multam libertatem in fide:ut intelligas gradum bonum in fi- de, ut crescant ex fide in fidem, qui ex fide bene mini- strant. Et multam libertatem in fide, id est abunde acci- piunt liberam in omnibus conscientiam, et fiduciam in deum, nulla cordis angustia et diffidentia laborantes, qui ex pura conscientia bene ministrant. quasi diceret Paulus, Tales eligendi funt in ministerium, ex quibus indies maior speratur per deum profectus.  \pend
\phantomsection
\addcontentsline{toc}{subsection}{\textit{In domo del. }}
\subsection*{\textit{In domo del. }}\pstart Haec domus deum habet habitatorem, regnum coel. rum, regnum dei: regnum Christi, templum fpiritus sancti, columna et stabilimentum ueritatis: quipp quam spiritus habitator fundauit supra firmam petra id est Christum, docens in ea omnem ueritatem: cou- tra quam ueritatem, ne portae quidem inferorum pr- ualere possunt: extra quam sunt merae tenebre, e- errores, et regnum principis tenebrarum, etc.  \pend
\phantomsection
\addcontentsline{toc}{subsection}{\textit{Vt noris quomodo, etc. }}
\subsection*{\textit{Vt noris quomodo, etc. }}
\phantomsection
\addcontentsline{toc}{subsection}{\textit{Magnum est pietatis mysterium }}
\subsection*{\textit{Magnum est pietatis mysterium }}\pstart Id est, arcanum, hactenus mundo incognitum, quoi ecclesia dei uiuentis nouit, cuius ministrisunt Timo- theus, Paulus etc. sicut ad Corin. dicit : Sic nos existi met homo ut ministros Christi, et dispensatores my- steriorum dei. Mysterium autem hoc est, deus manifi- status est in carne adsumpta. Ioan.I. Verbum carofi- ctum est, etc. et uidimus gloriam eius etc. Iustifica- tus est in spiritu:id quod deo non conuenit, sed Christi bumanitati, quam certum est non a seipsa fuisse iustam, sed a deo, qui solus iustificat. Ideo deo uero tribui deus et homo unus est Christus: sicut filius dei dicitux mortuus pro peccatis nostris. Recte ergo distribuit manifestatie  \pend\pstart Id est, quae curare debeas, etc. Nam mox capite se- quenti de nostris episcopis longe diuersis dicturus est  \pend
\section*{IN EPI. PAV. AD TIM. I. 89 }
\marginpar{[ p.89 ]}\pstart manifestatio est in carne, iustificatio autem inspiritu. In carne miraculis proditur, et praedicatur, in spiritu autem iustificatus, non nisi fide creditur. Siquidem in Christum hominem, deus ita iustitiam infudit, ut ipse esset nobis iustitia et sanctificatio et redemptio, sicut I. Corin.1. Qui factus est nobis a deo sapientia. et in psalmo 44. Propterea unxit te deus, etc. Et Ioan. o. Quem pater sanctificauit et misit in mundum. et Ioan.3. Non ad mensuram dat deus spiritum. Pater diligit filium, et omnia dedit in manum eius. Atque hic caro non prodest quicquam, nisi etiam iustificatum cre das Christum in spiritu, atque adeo iustitiam dei per quem iusti sumus, et satis fiat pro peccatis nostris et sal- uemur. Hoc autem mysterium manifestati dei in car- ne Christi, et iustificati dei in spiritu Christi, dignatio ne dei reuelatum est uniuersae creaturae. Nam et ange- li desyderant in ea prospicere, quae per Christum ha- bemus. 1.Petri 1. Et oculus non uidit, etc. Et praedi- tatum est gentibus, cum soli Iudaeisibi uiderentur ad Messiam pertinere, et ad mysterium regni coelorum, licut diximus ad Ephe. 2. et3. Creditum autem est in toto mundo: alioqui quid profuisset praedicatum esse? Praeterea receptum est in gloria, ut omnes glori ficent filium, sicut glorificant patrem:ut in nomine Ie lu omne genu flectatur, etc. et omnis lingua confite- atur quod Christus est in gloria dei patris. Gloria  \pend
\section*{ANNOT. IOANNIS POMERANI }\pstart haec est in omni terra, et superomnes coelos. Tante excellentior angelis factus est, Hebraeos I. Regnumo enim habet aeternum et sacerdotium, etc. Porro ce- tum est ueterem interpretem aliter legisse in Graeco sicut et Ambrosius legit et sic reddit, quod mani- festatum, etc. Et omnia quae sequuntur, ad mysterium referuntur:sed in eandemsententiam recidunt.  \pend
\endnumbering\beginnumbering\section{CAPVT IIII.}\pstart \huge\textbf{H}\normalsize Vnc locum explicaui in epistola secunda ad Thess. cap.2. Quae uero sint nouissima tem- pora, ut uetus translatio habet, dixi Esaiae 2.  \pend
\phantomsection
\addcontentsline{toc}{subsection}{\textit{Creauit ad sumendum. }}
\subsection*{\textit{Creauit ad sumendum. }}\pstart Gen.1. et 9. Creauitad sumendum cum gratiarum- actione, ut ex beneficio agnoscant patrem. Sanctus est qui comedit carnes, et gratias agit. Indignus est omni creatura, qui abstinet a carnibus, et non agno- scit creatorem. Lege 1. Corin. 1o. Quare libertas mea iudicatur, etc. Lege quoque ad Rho.14. Voluit et ta- li beneficio a gentibus agnosci, sed tenebrae erant ix mundo, sicut Paulus dicit Act.14. Qui in primis gens rationibus, et c. Sed quare dicit, fidelibus, quum et infi- deles edant, bibant, utantur creaturis? Dico, quod fi- delibus deus dedit suas creaturas, ut uides Gen. 1.ii est, ante lapsum Adae: post lapsum uero, infideles in ou mni peccant creatura. ad Titum I.Omnia munda mun  \pend
\section*{IN EPI. PAV. AD TIM. I. }
\marginpar{[ p.90 ]}\pstart uis, infidelibus nihil est mundum, etc. Hinc fit ut fide- lis edendo carnes, quacunque etiam die glorificet de- um. Hypocrita autem et iustitiarius abstinendo etiam peccat. Primum, quia damnat creaturam dei bonam, deinde gratias quidem simulare potest, agere non pont, quia abiicit dei beneficium. Et praeterea confidit in suam abstinentiam in contumeliam dei, cui soli fidendum est etc. Domini creati eramus omnium creaturarum, ue£ rum hoc dominium perdidimus in Adam, sed in Christo redditum est credentibus. Hinc est ille psal. Domini est terra, etc. Omnia subiecisti, etc. Ergo sine peccato utilis est omnis creatura dei, modo ne contra praeceptum dei. Nam traditiones humanas nihil moror, quia omnis creatura in Christo tua est: et ita quoque Paulus 1. ad Corin.1o. citat ex psal. testimontum. At dicunt iustiti- arii illi ciborum, et nos credimus. Respondet Paulus: mentimini. Ego fideles interpretor illos, qui cognoue runt hanc ueritatem, quam uos ignoratis, nempe, quod quica quid creauit deus, bonumsit, et nihil reiiciendum, si cum gratiarumactione comedatur.  \pend
\phantomsection
\addcontentsline{toc}{subsection}{\textit{Sanctificatur. }}
\subsection*{\textit{Sanctificatur. }}\pstart Quid est quod creatura dei perse bona est, et tamen hic dicit, sanctificatur? Respondeo bona est, quia cre- atura est: immunda autem est ei, qui ea abutitur, aut qui eam immundam ad usum credit. Sanctificatur igis tur ex usu utentis ea, sieut 1.Corin.7. dicitur. Sancti-  \pend
\section*{ANNOT. IOANNIS POMERANI }\pstart ficatus estuir, etc. Vbi sanctificatio ad usum matrimo nii pertinet, quod mulieri fideli licet uti matrimonio infidelis uiri, et contra absque peccato, et non minussine peccato quam si uterque esset fidelis. Hic non fit in se sanctus infidelis, sed ad usum matrimonii fit sanctus ut fidelis coniunx utatur infideli, ut re inse bona, et ad hoc a deo creata. Sic et cibus in se bonus, non fit sanctior, sed ad usum sanctificatur, ut sine peccatouta ris, quae infideli et gratias non agenti non esset san- cta, id est, munda. Sanctificatur, inquit, per uerbums dei, quo credimus omnia esse nobis sancta, ut supras et per orationem, qua talia oramus nobis dari, et pro datis gratias agimus, sicut supra dixit. Iamuide- stultitiam iustitiariorum, qui in cibis et aliis creatu- ris quaerunt sanctitatem, quum ipsae creaturae, licet bon- sint creatione, tamen sanctae nobis, id est mundae esse nem possunt sine nostra fide et gratiarum actione. Qua de usu ciborum diximus, omnia etiam intellige de usj matrimonii:id quod non obscure ex praecedentibus ins telligis. Creatura enim dei bona est mulier, si utaris sicut deus instituit. Si contra dei institutionem, ut fit in fornicatione et adulterio, iam bona tibi non est, quam deus bonam fecit. Iam et in eadem damnatione es, si in matrimonio putas usu creaturae (modo secundum deum fiat) peccari. Damnas enim non solum creaturam, sed et deiinstitutionem, etc.  \pend
\section*{IN EPI. PAV. AD TIM. I. }
\marginpar{[ p.91 ]}\pstart Iam doctrinae daemoniorum dicunt, Haereticus es.  \pend
\phantomsection
\addcontentsline{toc}{subsection}{\textit{Bonus eris minister. }}
\subsection*{\textit{Bonus eris minister. }}
\phantomsection
\addcontentsline{toc}{subsection}{\textit{Pietas ad omnia utilis, etCe }}
\subsection*{\textit{Pietas ad omnia utilis, etCe }}\pstart Pletas inscriptura (unde pis dicuntur )est vonor dei. Honoras autem deum si credis ei, fidis ei, times eum laudas, diligis, sciens omnia tua esse: in solo au- tem deo tuam esse salutem, et nullam rem, nullam Prorsus iustitiam humanam te posse liberare a peccas tis, aut saluare, sed solam dei gratiam per Christum Iesum dominum nostrum. Contra impietas est, ubi haec non adsunt, ut maxime omnia deceptis sensibus, et concupiscentiae oculorum sanctissima uideantur. Igi- tur exerceas te ad pietatem, non quod tuo studio ti bi talia parare nitaris, id quod fieri non potest: nam haec habuit iam Timotheus. Sed quando continuo de- loderio haec a deo postulas, sciens beatos esse qui esu- riunt et sitiunt iustitiam, ut tuam contine, uel sedulo lequi te uocationem. Et haec pietas ad omniautilis est:tua enim salus est et aliorum, et certa salus est, terta autem per uerbum dei. Nam promissiones ha- bet praesentis uitae et futurae, utraeque saluant, modo eis credas. Nam si uerbo dei credis, hinc iustificaris. Piis ut olim multa promittebantur terrena bona, ut est in ege Moysi, ita et nobis dicitur: Haec omnia adiicien- taruobis. Nam de promissionibus uitae futurae quid  \pend
\section*{ANNOT. IOANNIS POMERANI }\pstart attinet admontre? Corporalis autem exercitatio ta- les promissiones non habet, aliquid tamen utilitatis b- bet, scilicet dum non pro pietate eam habes, sed pro pietatis ministra: dum exerces corpus uigiliis, ieiuni- is, abstinentia, laboribus, ne insolescat contraspiritum et pietatem impediat: sed potius seruiat, sicut Christus admonet. Attendite uobis ne corda uestra grauentux crapula et ebrietate. Exercitatio enim est corporis quam exercere potes et intermittere pro tuo arbi- tratu, prout corporituo necessarium fuerit. Quod st corporali exercitatione, et quicquid externo opere fieri potest, te iustificari credas, iam corporalis exer- citatio non solum nihil habet utilitatis, sed etiam im- pia est, et te damnat: quia quod dei est tuis adscri- bis operibus, et adoras opus manuum tuarum, etc, Iacent igitur hic omnia iustitiariorum opera, et quae praeterea in fide salutis fieri non possunt, quia promil siones dei non habent. Longe ergo a salute omnium sectarum ordines. Atque hic addit Paulus obstruens con tradicentium ora. Indubitatus sermo etc. Hodie o- mnes fere de hoc dubitant, multi etiam blasphemant hanc ueritatem. Paulus indubitatum uult, et quod o- mnibus modis adprobetur: ut uero uniuersus mun- dus contemnet talia, ut maxime inter se omnes cons spirent contra nos, mundi sapientes et sanctuli illi attamen ab hac ueritate nos diuelli non sinamus,  \pend
\section*{IN EPI. PAV. AD TIM. I. }
\marginpar{[ p.92 ]}\pstart Illi in corporali exercitatione et ceremoniis sperats et mortuis rebus, quae nihil prodesse possunt, confi- dunt. Nos talia docemus esse impietatem, et spem fi£ xam habemus in deo uiuente, propterea non potest nos non odisse mundus. Haec confirment corda corum, qui nuper didicerunt pietatem solam esse, quam respiciat deus, quibus hic locus maxime commendatur.  \pend
\phantomsection
\addcontentsline{toc}{subsection}{\textit{Enutritus in sermone. }}
\subsection*{\textit{Enutritus in sermone. }}\pstart Duo dicit, oportet enim illum accepisse sermonem fidei, qui uult bene docere.  \pend
\phantomsection
\addcontentsline{toc}{subsection}{\textit{Praphanas et aniles. }}
\subsection*{\textit{Praphanas et aniles. }}\pstart Prophanum dici, quod sacrum non est, etiam puer nouit. Quicquid ergo ex sacris literis non est, propha num et anile est et fabula:licet hypocriiis et muu- dus hoc ipsum pro magna habeant sapientia, Christia nus homo hoc imo reponet pectori: quod si contempse rit, Christianus non est, etc.  \pend
\phantomsection
\addcontentsline{toc}{subsection}{\textit{Qui est saluator omnium, etc. }}
\subsection*{\textit{Qui est saluator omnium, etc. }}\pstart Certum est hoc dicere Paulum de saluatione tempo rali: alioqui non addidisset, maxime fidelium, quos sci- licet longe dignius saluat, quemadmodum et in psal- mo: Homincs et iumenta saluabis domine. Discant igitur gentes, discant hypocritae, non esse fidendum in creaturis et operibus suis, unde misere torquen- tur et pereunt, tantum abest ut salutem consequantur:  \pend
\section*{ANNOT. IOANNIS POMERANI }\pstart confidant autem in eo solo, unde etiam accipiunt cor, poris beneficia. Id quod negare non possunt.  \pend\pstart Si audis hic Paulum, eris haereticus hodie.  \pend\pstart Id est, sic uiue et exequere negotium tibi commis- sum, ut despici non possis: id quod protinus exponit dicens: Sed esto forma fidelium, etc.  \pend
\phantomsection
\addcontentsline{toc}{subsection}{\textit{Nemo tuam iuuentutem. }}
\subsection*{\textit{Nemo tuam iuuentutem. }}
\phantomsection
\addcontentsline{toc}{subsection}{\textit{Praecipe. }}
\subsection*{\textit{Praecipe. }}\pstart Id est, feruore spiritus, ut ad omnia sis strenuus et bonus miles Christi in fide, charitate, etc.  \pend
\phantomsection
\addcontentsline{toc}{subsection}{\textit{In spiritu. }}
\subsection*{\textit{In spiritu. }}\pstart Vult non esse neglectorem officii pastoralis, quod hic donum uocat. ad quod donum duo requiruntur. prophetia et uocatio. Sine his duobus nemo praesu- mat esse pastor uel episcopus. Prophetia est donum, ut possis potenter docere quae fidei sunt: quam si ha- bueris, expecta donec uoceris, uel a deo, uel ab homi- nibus, sicut omnes prophetae et apostoli uocatisunt. Vocatio tempore apostolorum fiebat per impositionem manuum super caput uocati, quo ceu signo agnoscebae tur ille uocatus et susceptus. Sic non apostoli (ut non sibi hoc arrogent nostri apostoli) sed discipuli alii ab apostolis congregati eligunt septem diaconos, et orantes  \pend
\phantomsection
\addcontentsline{toc}{subsection}{\textit{Donec uenero, etc. }}
\subsection*{\textit{Donec uenero, etc. }}
\section*{IN EPI. PAV. AD TIM. I. 93 }
\marginpar{[ p.93 ]}\pstart imponunt eis manus in conspectu apostolorum: bo- signo ostendentes apostolis, et reliquae multitudini, quod tales uolebant, et elegerunt ministros. et Paulus imposuerat manus Timotheo docto sacris literis a puero, et erudito in fide, sicut dicit epistola secunda cap.1. Sed non solus imposuerat, ut hic uides. Nam quod hic translatum est, authoritate sacerdotii, non est in Graeco: sed μετα ἐπιθεσέωσ, etc. id est, cum impo- sitione manuum presbyteri: quasi dicas latine, senatus, id est seniorum. Vides consensu aliorum Timotheum ele- ctum: ita et intellige quod infra dicitur, manus cito nemini imponas. Et quod Tito scribitur, ut constituas per ciuitates presbyteros, etc. Impositio manuum fuit Hebraeis familiaris, ubi aliquid commendabatur deo. Sic Iacob imponit manus duobus filiis Ioseph, Genes. 48. Sic manus super caput hostiae immolandae ponitur, Leuitici I. Sic Christus imponit manus super paruu- los: et de credentibus dicit, Super aegros manus impo nent, et bene habebunt. Hoc quia externum signum est, uel omitti potest, uel aliter fieri: quia institutum non est, aut mandatum. Electionis tantum signum erat. Quae institutasunt, ut quod oportet episcopumesse sine crimine, aptum ad docendum, etc. Hodie et negli gimus, et contemnimus, et uana signa magnifacimus: uana dico, ubi res non adest. Et insupersignis externis omnia complemus, et fingimus talia, quae Christus et  \pend
\section*{ANNOT. IOANNIS POMERANI }\pstart apostoli Christi non nouerunt.  \pend
\phantomsection
\addcontentsline{toc}{subsection}{\textit{Haec exerce. }}
\subsection*{\textit{Haec exerce. }}\pstart Vt supra, admonet eum sui officii. Haec exerce, id est in his persiste.  \pend
\endnumbering\beginnumbering\section{CAPVT V}
\phantomsection
\addcontentsline{toc}{subsection}{\textit{Seniorem ne increpaueris, etc. }}
\subsection*{\textit{Seniorem ne increpaueris, etc. }}\pstart \huge\textbf{H}\normalsize Aec non est adulatio. Nam post dicitur in hos ipso cap. Peccantes coram omnibus, argue, ut et caeteri timorem habeant. Sed potius est Christiana modestia cum omnibus, tum maxime Christiano docto- ri necessaria, ut et uerba quoque nostra spirent amo- rem. Verum hic additur, cum omni castitate, ut blanda foeminarum appellatio non solum apud nos pura, sed etia apud ipsas et apud alias sine suspicione sit. Oportet enim nos prouidere bona, non tantum coram deo, sed etiam coram hominibus, ad Rhom.12.  \pend
\phantomsection
\addcontentsline{toc}{subsection}{\textit{Viduas honora. }}
\subsection*{\textit{Viduas honora. }}\pstart Locus hic est de uiduis, quas omni auxilio destitu- tas, suscipiebat ecclesia alendas, quibus prouidebant Diaconi ex bonis ab ecclesia in pauperes collatis. Hae non propter uotum, quod scriptura non nouit uspiam suscipiebantur, sed propter paupertatem. Et has uo- cat uere uiduas, quae non solum marito, sed etiam filiis aut amicis, qui eam alere possent, aut bonis destitutae  \pend
\section*{94 IN EPI. PAV. AD TIM. I. }
\marginpar{[ p.94 ]}\pstart funt, has praecipit honorandas, non uerbis aut gestu: sed rebus ipsis, sicut praecipiuntur honorari parentes. et infra: Presbyteri duplici honore digni sunt, etc. Vide quaeso quam diligenter cauet Paulus ne iniiciat laqueum, aut certe ne cadatur in laqueum diaboli ab ipsis uiduis, quae tantum suscipiebantur alendae, ne ocium sectantes in alieno labore inciperent lasciui- entes deficere a Christi fide. Quam bene nouit Pau- lus hoc genus, de quo dicit, quae in deliciis uersatur, ea uiuens quidem carne, mortua tamen est deo: id est, fidem abnegauit. Vbi autem sunt illi laudatores uir- ginitatis et uiduitatis? qui nihil aliud faciunt quam quod in laqueos trudunt insanis laudibus, nescientes de quibus adfirmant. Paulus laqueos cauet, et in uir- ginitate et uiduitate, sicut hic uides. et I. Corin.7. Propter fornicationem, inquit, unusquisque habeat uxoremsuam. Et iuniores uiduas nubere uolo, filios procreare, etc. Et sic sequentur non Hieronymum, led deum, qui eas ad talia creauit, certae quod deo ta- lia placeant, et quod per haec seruiant proximo, id est, marito, filiis, familiae. Reliquain hunc locum uide in libel lo de uotis in fine. De uiduis hodie minus periculi est, de quibus tamen solis scripsit Paulus, sed alendis, non uo to adstringendis. Doctrinae uero daemoniorum perpe- tuos laqueos omni fere aetati et sexui iniecerunt, etc.  \pend\pstart Quod si qua uidua.  \pend
\section*{ANNOT. IOANNIS POMERANI }\pstart Viduae annos natae sexaginta, aut sunt uere uiduae et desolatae, hae alantur: aut non sunt desolatae, sed habent liberos aut nepotes, hae gubernent domum pie, id est, in domino et doctrina fidei, et uicem reddant maio- ribus: id est, sic seruiant minoribus ad se pertinentibus (quos uocant liberos et nepotes) sicut ipsis seruitum est a suis maioribus. Quod si facere nolunt, fidem ab- negarunt, et sunt infidelibus deteriores. Quam enim fidem haberent qui hoc nolunt quod deus uult, et illos abiiciunt, quos deus ipsis dedit, quibus nisi seruiant, nulli seruire possunt? Qua quaeso charitate uterentur in alienos, quae suos negligunt et contemnunt? Fides esse non potest, ubi fructus eius non ostenditur, nonso- lum erga alienos, sed etiam erga illos ad quos etiam na- tura propensa est.  \pend
\phantomsection
\addcontentsline{toc}{subsection}{\textit{Porro quae uere uidua est. }}
\subsection*{\textit{Porro quae uere uidua est. }}\pstart Describit quae sit uere uidua, quae desolata speratur caste uictura, tamen si iunior est, nubat et non suscipia- tur alenda, tradatur marito alenda, ne postea lasciuiat in Christum. Sed quis hodie pauperem tradit marito? Felix haec ordinatio, sed non seruata. Degenerauit autem in sectas perditionis utriusque sexus.  \pend
\phantomsection
\addcontentsline{toc}{subsection}{\textit{Hominum testimonio, etc. }}
\subsection*{\textit{Hominum testimonio, etc. }}\pstart Nec omnes uiduas sexaginta annorum, licet desolas taesint, suscipit: sed quae testimonio hominum compro- batae sunt, ne postea sint scandalo Christianis inter  \pend
\section*{IN EPI. PAV. AD TIM. I. }
\marginpar{[ p.95 ]}\pstart gentes. Irreprehensibiles, id est, sine crimine coraus bominibus.  \pend\pstart Id est, fidem in Christum, alioqui non diceret, lasciuire toeperint in Christum: quae deinde non solum sunt scan dalo Christianis, uerumetiam obloquuntur in domi- bus aliorum, ctiam gentilium, dum quaerunt maritos. Vide libellum de Votis. Hoc est quod hic addit, nullam occasionem dare aduersario, ut habeat ansam maledi- cendi Euangelicae doctrinae. Fidem uero in Christum et non uotum intelligere Paulum satis ostendit quod Iequitur. Iam enim nonnullae deflexerunt, secutae satanam.  \pend
\phantomsection
\addcontentsline{toc}{subsection}{\textit{Primam fidem. }}
\subsection*{\textit{Primam fidem. }}
\phantomsection
\addcontentsline{toc}{subsection}{\textit{Quod si quis fidelis. }}
\subsection*{\textit{Quod si quis fidelis. }}\pstart Iam etiam gloriantur monachi, quod sua relinquant, at uiuant ex ecclesiae bonis. Clerici uero et sua retinent, et pauperum sudore uiuunt. Paulus non uult, ne uiduis quidem, quibus per amicos succurri potest, grauari ecclesiam.  \pend\pstart Hactenus de uiduis, etc. Praeceptum charitatis exis git uere uiduis succurrere. Caetera quae hic ordinat de uiduis Paulus, bona sunt ad cauendos multos laqueot in uiduis: quae tamen non ita accipe, quasi legem tibi praescribat Paulus, cum uir spiritualis posset etiam adhuc aliqua ex his immutare, et aliter obseruare pro loco et personis, si uideret expedire. Praeterea de quibus scribit, hoc honestum est et acceptum coram  \pend
\section*{ANNOT. IOANNIS POMERANI }\pstart deo. Tu uero uide ne ista contemnas, caue laqueos, caue scandala. Paulus astutias satanae non ignorauit, etc.  \pend
\phantomsection
\addcontentsline{toc}{subsection}{\textit{Qui bene praesunt presbyteri. }}
\subsection*{\textit{Qui bene praesunt presbyteri. }}\pstart Senioribus, qui praesunt in ciuitate reliquae multitu dini, et bene praesunt, debetur duplex honor. Alter, ut authoritati eorum deferatur, et obediatur eis cum reuerentia. Bene enim praesunt, id est secundum iudi- cium spiritus quod habent, omnia sine neglectu dispo- nunt. Alter, ut ex ecclesiae bonis eis prouideatur de necessariis, qui ob gubernandi et curandi quae greg salubria sunt, occupationem et studium, sua curare non possunt. Et cum haec debeantur illis, qui bene praesunt. id est in domino praesunt, in primis debentur etiam his, qui laborant in sermone et doctrina, id est euan- gelistis siue doctoribus, Matthaei Io. Iad Corin.9. Gala.6.  \pend
\phantomsection
\addcontentsline{toc}{subsection}{\textit{Aduersus presbyterum. }}
\subsection*{\textit{Aduersus presbyterum. }}\pstart Qui praesunt, facile incurrunt odia inobedientium, et idcirco non facile audiri debent, qui tales accusant, cum quod suspecti sunt accusatores, tum quod maiori scandalo sunt, qui praesunt, si crimine notati fuerint uel suspecti. Quod si manifestum fuerit ecclesiae eos peccal se, aut peccare, uel contra eos probatum fuerit duobus aut tribus testibus secundum legem, Deu.19. Ab epi- scopo, id est euangelico doctore siue apostolo, qui illic constitutus est, debent coram omnibus argui, ne dum  \pend
\section*{IN EPI. PAV. AD TIM. I. }
\marginpar{[ p.96 ]}\pstart conniuetur, eorum error uideatur uel approbari, uel non esse error, et inueniat imitatores. Ideoque non est necessarium quaerere, quomodo haec conueniant cum uerbis Christi, Matt. 18. Corripe eum inter te et ipsum. Nam hic dictum est ecclesiae. Ergo coram ecclesia cora ripiendi sunt presbyteri peccantes, si sic conuicti fue- rint. Maius praeterea periculum est in his qui prae- sunt, si peccant, quam in aliis fratribus. Aduerte er- go, quod consequentia textus postulat, ut quod dicit, eos qui peccant, etc. de senioribus siue presbyteris in telligas. Si uero de quibuscumque intelligere uolueris, ta men necesse est, ut de manifestis criminibus, aut de con uictis intelligas secundum euangelium.  \pend
\phantomsection
\addcontentsline{toc}{subsection}{\textit{Obtestor in conspectu dei, etc. }}
\subsection*{\textit{Obtestor in conspectu dei, etc. }}\pstart Deus, Christus, et electi angeli nostra uident: sic ergo exequamur nobis commissa, ut in conspectu ta- lium. Vt serues, inquit, haec sine praecipitatione iudicii, sed sano consilio et diligenti inquisitione per testes, etc. Nihil facies iuxta propensionem animi, id est car nali iudicio et fauore in persona, aut odio, hoc est, uide ne subito proferas sententiam secundum affectum subitum, quia et lex dei de iudicibus dicit, Deute.16. Iudicent populum iusto iudicio, nec in alteram partem declinent, non accipias personas, nec munera: quia mu nera excaecant oculos sapientium, et mutant uerba  \pend
\section*{ANNOT. IOANNIS POMERANI }\pstart iustorum, iuste quod iustum est prosequaris.  \pend\pstart Id est subito, ne alicui committas episcopi munus, sel probetur antea, si talis sit, qualem supra depinxi.  \pend
\phantomsection
\addcontentsline{toc}{subsection}{\textit{Neque communices peccatis alienis. }}
\subsection*{\textit{Neque communices peccatis alienis. }}
\phantomsection
\addcontentsline{toc}{subsection}{\textit{Manus cito, etc. }}
\subsection*{\textit{Manus cito, etc. }}\pstart Si qui iniuste faciunt, si qui quod iustum est, non iu- ste exequuntur, uide ne consentias, ne conniueas, ne taceas: sed contradicas, etc. Vt si nihil aliud efficias, tamen conscientiam tuam liberes, hoc est quod addit; temetipsum purum serua.  \pend
\phantomsection
\addcontentsline{toc}{subsection}{\textit{Ne posthac bibas aquam. }}
\subsection*{\textit{Ne posthac bibas aquam. }}\pstart Hinc habes, quod corporales exercitationes pro nostro arbitrio seruare, mutare, intermittere, omitte- re possumus, quas adhibemus secundum corporis no- stri rationcm, ut seruiat, non intereat, etc.  \pend
\phantomsection
\addcontentsline{toc}{subsection}{\textit{Quoruudam hominum peccata. }}
\subsection*{\textit{Quoruudam hominum peccata. }}\pstart Quia dixerat, neque communices peccatis alienis:ne hic superstitiosius trepidaret Timotheus, dicens: Quid si ignorem peccata aliena, id est illorum quibus com- munico? respondet: Dies docebit. Ergo quod dixi, ne sis in iudicio praeceps, nihil opertum, quod non reue- labitur: interim nihil iudicato. Pro bonis habeto, de quibus nemo male sentit, qui nemini sunt occasio pec- cati. Hic nihil tuae conscientiae, aut tuo officio erit pe- riculi. Atque hanc hic esse ueram sententiam, quam con  \pend
\section*{IN EPI. PAV. AD TIM. I. }
\marginpar{[ p.97 ]}\pstart textus requirat: credo, ut iudicium hic non intelligas iltud Christi extremum, sed iudicium humanum in hac£ uita, quod supra commisit Timotheo. Sic dicit: Mani- festa peccata procedunt ad iudicium, id est non indi- gent iudicio, uel ut iudices, siue sententiam de eis pro- feras, quia iam iudicata sunt. Aliquorum uero peccata lubsequuntur: id est, postea nota erunt. Ex fructibus enim cognoscetis eos. Sic et opera bona manifesta sunt ante: id est, antequam iudicentunuel de eis sententia feratur. Ea uero quae secus habent, id est in occulto ma- le fiunt, aut uelut bona sese per hypocrisin uenditant, oc cultari non possunt, id est suo tempore prodentur. Tunc ergo iudicandi sunt ab ecclesia mali, quando malos es- Ie notum fiet. Non maculauit Christum, quod Iudam ha- buit apostolum, caeteris nescientibus, quod esset diabolus et proditor. Nam interim non male docuit, officium non neglexit, nemini scandalo fuit, et c. Sic et de manifestis Fuctibus carnis et spiritus legis ad Gal.5.  \pend
\endnumbering\beginnumbering\section{CAPVT VI}
\phantomsection
\addcontentsline{toc}{subsection}{\textit{Quicunque sub iugo sunt serui, etc. }}
\subsection*{\textit{Quicunque sub iugo sunt serui, etc. }}\pstart \huge\textbf{P}\normalsize Rima ad Corin.7. dicitur: Vnusquisque in qua uo catione uocatus est frater, in hac permaneat a- pud deum. Apud deum, inquit, ne credas te in conditios ne turpi oportere manere, quando uenisad fidem. De oibus ergo mediis intellige. Media sunt siue indiffe-  \pend
\section*{ANNOT. IOANNIS POMERANI }\pstart rentia omnia externa, praeter ed quae prohibet. Non ergo refert, Christianus siue Iudaeus sit, siue gentilis carne, siue seruus sit, siue dominus, siue mas sit siue foe- mina, siue sutor siue consul, siue rusticus siue princeps, etc. Serui igitur non negligunt officium, quia Christi- ani facti. Nam hanc conditionem eis deus dedit, ut in bac uitam transigant, nisi melior sese libertatis etiam corporalis obtulerit occasio, sine domini iniuria, aut cum domini uoluntate. Seruient ergo cum bona consci- entia: siue dominis gentilibus, quia hoc gratum deo est, ne blasphemetur Euangelica grat ia, sed commende tur potius, quod tales faciat seruos, talesque homines, siue Christianis dominis: quia iam non tam dominis, quam fratribus seruiant, de quo etiam gaudeant, quod tales ha beant dominos, qui turpia exigere non possunt. De hoc loco uide in epistola Petri, et saepe alias in Paulo.  \pend\pstart Omnia repetit quae dixit. Nam (quod et in princi- pio monuit (aliud docentes uitari uult: id est, illos (ut ait) qui non accedunt sanis sermonibus Christi, et do- ctrinae quae ad pietatem tantum pertinent. Hic uides haereticos esse omnes, et excommunicandos, qui aliud docent. Vides quoq omnes epistolas Pauli nihil aliuq esse, quam sermones Christi, et doctrinam quae so- lum spectet pietatem. Quomodo enim alia doceret, quae ipse excommunicationis maledicto uetat? Aduer.  \pend
\phantomsection
\addcontentsline{toc}{subsection}{\textit{Haec doce, etc. }}
\subsection*{\textit{Haec doce, etc. }}
\section*{IN EPI. PAV. AD TIM. I. 98 }
\marginpar{[ p.98 ]}\pstart te hic quam egregie depingat illud genus doctorum, qui aliud docent, scilicet esse homines inflatos, ignoran tes, insanientes ad inutiles quaestiones et disputatio- nes de genealogiis, de iudaicis fabulis, de animabus se- paratis, de choris angelorum, de fictis reuelationibus, de humanis traditionibus: de his quae scriptura non no- uit, ex quibus nascitur inuidia, etc. dum quisque se prae- ferre sua sapientia nititur. In summa dicit esse homi- nes, quibus adempta est ueritas: id est, fides in Christum, qui omnia tentant, qui quaestui suo consulant:uocant tamem interim quaestum pietatem: id est, gloriam et cul- tum dei. Hic uides so phisticos theologos, et papisticum regnum, de quibus rursum in fine. O Timothee, etc. ubi iubet eos excommunicandos, et eorum doctrinam uitandam, quod habeant ficta et inania uocabula, quae seriptura non nouit, et opponant scientiam suam quae scientia non meretur dici, sed inscitia et insipientia: quam pfitentes, a fide et a scopo uero aberrauerunt. Isti sunt ꝗ liberum arbitrium habent pro gratia, rationem Pfide, opiniones pro iudicio spiritus, doctrinas homi- num pro doctrina sacra:qbus nihil in ore est, quamipri- matus Petri, fides acquisita:uniuersalia, formalitates, actus elicitus, intellectus agens, intellectus passibilis, comnotata, relationes rationis in peccatis, etc. ꝗ quasi doctissimi quauis de re disputant, et aliter sentientibus Ie sua scientia opponunt, etc. Sed hoc est, inquit Paulus.  \pend
\section*{ANNOT. IOANNIS POMERANI }\pstart a fide longe abesse, et non esse Christianos, qui se sa- pientissimos, et haereticae prauitatis extirpatores ia- ctitant. Seiungere, inquit, ab his qui eiusmodisunt.  \pend\pstart Quid attinet falsa scientia et humana doctrina pro quaestu seducere mundum, cum quaestus uerus sit et magnus, habere pietatem, et esse sua sorte contentums O uerbum breue et suauissimum, quod creditum fa- cit in utraque uita beatos, et mente pacatos, Contra qui non curata pietate, sua sorte nolunt esse contenti sed uolunt ditescere: id est, uoluntatem quam deo de- bent, consecrant mammonae, incidunt in tentationem et laqueum, etc. quia pro consequendo lucro incipi- lint agere contra conscientiam, et perdunt fidem, sis eut uides in Euangelio de semine, quod cecidit inter spinas. Caeterum attende quod dicit, qui uolunt diuites fieri, non contenti sua sorte: et non dicit, qui diuites sunt. Nam diuitiae sunt media, qbus spiritualis recte uti potest. Hoc est quod infra dicit: his qui diuites sunt etc. ubi praescribit quod sit recte uti diuitiis, si qui eas habent,  \pend
\phantomsection
\addcontentsline{toc}{subsection}{\textit{Est auten quaestus, etc. }}
\subsection*{\textit{Est auten quaestus, etc. }}
\phantomsection
\addcontentsline{toc}{subsection}{\textit{De laboribus multis. }}
\subsection*{\textit{De laboribus multis. }}\pstart Id cum accidat omnibus qui solicitudinem deo noù relinquunt secundum Euangelium: Nolite soliciti es- se, etc. Quid putas futurum de illis, qui contra conscienti- am et pietatem audent quaerere sua lucra? etc,  \pend
\section*{IN EPI. PAV. AD TIM. I. }
\marginpar{[ p.99 ]}\pstart Nam ad eum non pertinet, qui uult seruire mammonae.  \pend
\phantomsection
\addcontentsline{toc}{subsection}{\textit{Tu uero homo dei. }}
\subsection*{\textit{Tu uero homo dei. }}
\phantomsection
\addcontentsline{toc}{subsection}{\textit{Sectare iustitiam, etC. }}
\subsection*{\textit{Sectare iustitiam, etC. }}\pstart Per haec cum deo agimus. Sectare charitatem, patien tiam, mansuetudinem: per haec cum proximo agimus. Hae sunt diuitiae Christianorum, de quibus Christus dicit: Primum quaerite regnum, etc.  \pend
\phantomsection
\addcontentsline{toc}{subsection}{\textit{Certa bonum certamen fidei. }}
\subsection*{\textit{Certa bonum certamen fidei. }}\pstart Persta in fide, et nullis cede uel aduersis, uel aduer ariis, ut agas uel omittas contra tuae rationem fidei: id est, contra id quod iam doctus sentis. Adprehende uitam aeternam, ut hanc uitam pro ueritate contemnas, ne huius uitae amore contra fidem aliquid committas, aut in aliquo cedas aduersariis. Ad quam, inquit, uocas tus es. Christianus homo ad mortem huius uitae, et quicquid ex Adam nostrum est, uocatur, et non perficitur eius sanctificatio, donec omnia intereant. Tantum e- nim apprehendit uitae aeternae, quantum ei quotidie baptizante spiritu et igne, decedit ex hacuita. Haec est mortificatio ad uitam. Dum ergo uocaris ad uitam, certus esto hoc ipsumsine morte fieri non posse. Sine cruce non glorificatus est Christus, ita nec Christia- nus sine cruce glorificabitur.  \pend
\phantomsection
\addcontentsline{toc}{subsection}{\textit{Et professus es bonam professionem. }}
\subsection*{\textit{Et professus es bonam professionem. }}\pstart Scilicet in baptismo, cui solebant Christiani adesse,  \pend
\section*{ANNOT. IOANNIS POMERANI }\pstart ut sciretur hic esse frater qui baptizabatur, cuius iam confessionem adstantes, ceu testes audirent. Nisi hic uee lis intelligere Timotheum hanc professionem docene- di populum, et episcopi munus suscepisse coram mul- tis testibus: id est, praesente populo Christiano.  \pend
\phantomsection
\addcontentsline{toc}{subsection}{\textit{Praecipio tibi coram ideo, etc. }}
\subsection*{\textit{Praecipio tibi coram ideo, etc. }}\pstart Vides in hac epistola, quanta cura tommendet Ti- motheo, quae sunt muneris episcopi: praecipit et adiu rat per omnia Christianorum sacra, sicut saepe uides: id quod non facit in aliis epistolis. Nimirum ostendens in hoc negocio pendere omnia: quod si non pie cure- tur, relicta doctrina pietatis, oia quae Christiana sunt intereunt, sicut factum uidemus. Vnde hic praecipit coram deo, qui uiuificat omnia: et coram Christo, qui a pro- fessione sua non decessit, cum erat adiudicandus mor- ti, ut uitam ad quam uocatus est Timotheus, et pro- fessionem quam professus est, mordicus teneat, etiam inconspectu mortis et inferorum:id est, ut seruet man datum quod accepit, immaculatus coram deo et irre prehensibilis coram hominibus usque in finem. Imprimis enim necessarium est episcopo, ut faciat quae sui muneris funt: ad quod certe indiget non paruo spiritu, necesse que est, ut iugiter postulet a deo: alioqui nihil efficiet,  \pend
\phantomsection
\addcontentsline{toc}{subsection}{\textit{Vsque ad apparitionem. }}
\subsection*{\textit{Vsque ad apparitionem. }}\pstart Christiani ueri sic omnia deponunt, ut parati et uis  \pend
\section*{IN EPI. PAV. AD TIM. I. }\pstart gilantes possint occurrere Christo uenienti ad iudici- um, scientes et se tunc glorificandos. Sicut de quinque  prudentibus uirginibus habet parabola. Id quod legis fere ubique in nouo testamento, praecipue in concioni- bus Iesu saluatoris, et epistolis Pauli:ad quem diem hac uita perfuncti omnes seruantur exuscitandi, et ta les seruantur, quales hinc abierimus. Tunc omnes si- mul reprobi damnabuntur, omnes simul electi glori- ficabuntur.  \pend
\phantomsection
\addcontentsline{toc}{subsection}{\textit{Et solus habet immortalitatem. }}
\subsection*{\textit{Et solus habet immortalitatem. }}\pstart Frustra ergo aliunde petieris, ut sis immortalis. Ita Jolus habet immortalitatem, sicut solus habet uitam, qui et ipsa uita est. Caeterum de immortali et beato deo dixi quoque supra cap.I.  \pend
\phantomsection
\addcontentsline{toc}{subsection}{\textit{Lucem inaccessam. }}
\subsection*{\textit{Lucem inaccessam. }}\pstart Id est, hominibus incomprehensam: id quod exponit dicens, quem uidit nemo, etc. Vide annotata Ioan.I. et Esaiae 6. Vnde ergo cognoscimus deum, sine quo non possumus esse beati? Respondeo, lege I. Ioan.4. et inuenies.  \pend
\phantomsection
\addcontentsline{toc}{subsection}{\textit{Depositum. }}
\subsection*{\textit{Depositum. }}\pstart Id quod tibi creditum est, scilicet doctrinam et oh cium.  \pend\pstart FINIS  \pend
\endnumbering
\end{pages}
\end{document}
        